6. Conclusion


This paper provides the first systematic, reproducible description of Ecuador's homicide wave at the provincial level using exclusively official data, and a transparent assessment of the labor market channel invoked in public debate. Between 2018 and 2025, intentional homicides rose from 996 to 9,283 -- a rate of approximately 51 per 100,000 -- with 4,154 additional deaths in the first half of 2026. The wave accelerated from 2021 (2,495) through 2023 (8,248), is national in scope despite Guayas' 44.2\% share of accumulated homicides, and preceded the "internal armed conflict" decrees of 2024--2026. Provincial labor conditions, measured from the single available ENEMDU cross-section (2024-II: unemployment 3.48\%, underemployment 20.97\%, adequate employment 35.01\%), are weakly associated with homicides: Pearson correlations do not exceed 0.34, the Spearman correlation for underemployment in 2025 reaches 0.42, and regression coefficients are small and statistically insignificant ($\beta$ = +0.08, p $\approx$ 0.39 for unemployment; $\beta$ = $-$0.04, p $\approx$ 0.47 for underemployment). The event study around Decreto 111 detects no significant differential change in homicides, the placebo is flat at the pseudo-event, and the Esmeraldas synthetic control is suggestive only.


These results align with the verified international literature: unemployment robustly predicts property crime, but homicide waves in Latin America are explained by illicit markets, armed-group fragmentation, and state capacity, not by aggregate unemployment (Chiricos, 1987; Fajnzylber et al., 2002a; Dell, 2015; Castillo et al., 2020; Lessing, 2021). The implication is not that labor markets are irrelevant -- precariousness lowers the opportunity cost of illegal economies for young people, and gang territorial control destroys licit employment (Melnikov et al., 2025; Sviatschi, 2022a) -- but that the labor channel operates through the quality of licit alternatives, adjusts slowly, and cannot explain the timing or magnitude of the explosion.


The paper's central methodological lesson is that the current data cannot support causal inference on the labor--crime nexus in Ecuador: a single ENEMDU cross-section absorbs the labor signal in province fixed effects; 24 provinces make conventional clustered inference fragile; and no credible instrument or eligibility threshold exists. These are design gaps, not omissions. A quarterly provincial ENEMDU panel since 2018, canton-level data, interoperable administrative records, and -- where treatment dates vary across provinces -- modern staggered estimators would transform the menu of identifiable questions, as would double/debiased machine learning on individual microdata.


For policymakers, the findings argue for an evidence-based division of labor: focused policing and state presence in high-lethality territories with measurable targets and civilian oversight; evaluation of the Bono de Desarrollo Humano and youth employment programs with causal standards and labor-insertion rather than homicide targets; and a technical evidence unit empowered to evaluate the exceptional regime and publish results regardless of sign. None of these recommendations promises causal effects that the evidence cannot support; institutionalizing that honesty is the paper's most direct policy contribution.


Finally, the paper is an exercise in stated limitations: every quantitative claim has been checked against its source, every causal claim explicitly withdrawn, and every policy recommendation labeled by the strength of the evidence behind it. Ecuador is living through a public-health and governance emergency whose drivers lie primarily outside the labor market. Better data, better evaluation, and disciplined public discourse about what is known -- and what is not -- are the preconditions for a security policy that can be held accountable.